% Electroencephalogram (EEG) is a brain signal known for its high time resolution but moderate signal-to-noise ratio. Whether natural images can be decoded from EEG has been a hot issue recently. 
Electroencephalography (EEG) signals, known for convenient non-invasive acquisition but low signal-to-noise ratio, have recently gained substantial attention due to the potential to decode natural images. 
This paper presents a self-supervised framework to demonstrate the feasibility of learning image representations from EEG signals, particularly for object recognition. 
The framework utilizes image and EEG encoders to extract features from paired image stimuli and EEG responses. 
Contrastive learning aligns these two modalities by constraining their similarity.
% With the framework, we attain significantly above-chance results on a comprehensive EEG-image dataset, achieving a top-1 accuracy of 15.6\% and a top-5 accuracy of 42.8\% in challenging 200-way zero-shot tasks. 
Our approach achieves state-of-the-art results on a comprehensive EEG-image dataset, with a top-1 accuracy of 15.6\% and a top-5 accuracy of 42.8\% in 200-way zero-shot tasks.
Moreover, we perform extensive experiments to explore the biological plausibility by resolving the temporal, spatial, spectral, and semantic aspects of EEG signals. 
Besides, we introduce attention modules to capture spatial correlations, providing implicit evidence of the brain activity perceived from EEG data.
These findings yield valuable insights for neural decoding and brain-computer interfaces in real-world scenarios.
%Electroencephalogram (EEG) is a brain signal known for its high time resolution and moderate signal-to-noise ratio. 
% Whether natural images can be decoded from EEG has been a hot issue recently.
% In this paper, we propose a self-supervised framework to demonstrate the feasibility of learning image representation from EEG signals with the object recognition task.
% Specifically, image and EEG encoders are first used to extract features from paired image stimuli and EEG responses. 
% Then, we employ contrastive
% learning to align these two modalities by constraining their similarity.
% We introduce two plug-in-play modules that capture spatial correlations for better features.
% Our approach achieves state-of-the-art results on a large and rich EEG-image dataset, with a top-1 accuracy of 15.6\% and a top-5 accuracy of 42.8\% in 200-way zero-shot tasks.
% Extensive experiments demonstrate good biological plausibility by resolving the temporal, spatial, spectral, and semantic aspects of EEG signals.
% These results offer valuable insights for neural decoding and real-world applications of brain-computer interfaces.
Code available at https://github.com/eeyhsong/NICE-EEG.

% ！重构三个contribution
% 1. 提出了自监督模型应用于这个场景，取得了很好的效果；
% 2. 证明可行性，自然图像，获得语义信息；
% 3. 空间滤波器，给出模型关注有效特征的证据，进一步帮助